\documentclass[a4paper,10pt,fleqn]{article}
\usepackage[margin=1in]{geometry}
\usepackage{amssymb}
\usepackage{adjustbox}
\usepackage{natbib}
\usepackage[colorlinks=true,linkcolor=blue,citecolor=blue,urlcolor=blue]{hyperref}
\usepackage{fuzz}
\usepackage{schemapk}
\usepackage{zed-maths}
\usepackage{zed-proof}
\newdimen\savedleftskip
\begin{document}

\section*{Syntax Environment Examples}
\addcontentsline{toc}{section}{Syntax Environment Examples}

\begin{zed}
[NAME]
\end{zed}

\subsection*{Example 1 : Simple Enumeration}
\addcontentsline{toc}{subsection}{Example 1 : Simple Enumeration}

\noindent For simple free types with several branches, syntax provides column-aligned layout:

\bigskip

\begin{syntax}
OP1 & ::= & plus | minus | times | divide
\end{syntax}

\noindent This generates aligned columns with the ::= separator clearly visible.

\bigskip

\subsection*{Example 2 : Parameterized Constructors}
\addcontentsline{toc}{subsection}{Example 2 : Parameterized Constructors}

\noindent Constructors can have type parameters using angle brackets:

\bigskip

\begin{syntax}
Tree & ::= & leaf \ldata \nat \rdata | branch \ldata Tree \cross Tree \rdata
\end{syntax}

\noindent Parameters are rendered with special data delimiters in LaTeX.

\bigskip

\subsection*{Example 3 : Multiple Definitions with Grouping}
\addcontentsline{toc}{subsection}{Example 3 : Multiple Definitions with Grouping}

\noindent Blank lines between definitions create visual groups using \textbackslash also:

\bigskip

\begin{syntax}
Status2 & ::= & active2 | inactive2 | pending2 \\
Response & ::= & ok | error \ldata \nat \rdata | timeout
\end{syntax}

\noindent The blank line causes the also command to be generated, creating vertical separation between the two type definitions.

\bigskip

\subsection*{Example 4 : Complex BNF - Style Definitions}
\addcontentsline{toc}{subsection}{Example 4 : Complex BNF - Style Definitions}

\noindent Syntax environment is ideal for BNF-style grammar definitions:

\bigskip

\begin{syntax}
Expr & ::= & const \ldata \nat \rdata | var \ldata NAME \rdata | binop \ldata OP1 \cross Expr \cross Expr \rdata \\
Stmt & ::= & assign \ldata NAME \cross Expr \rdata | seq \ldata Stmt \cross Stmt \rdata | skip
\end{syntax}

\noindent This creates a professionally formatted grammar specification.

\bigskip

\subsection*{Comparison : Syntax vs Zed Blocks}
\addcontentsline{toc}{subsection}{Comparison : Syntax vs Zed Blocks}

\noindent Standard zed blocks work for simple types:

\bigskip

\begin{zed}
Color1 ::= red1 | green1 | blue1
\end{zed}

\noindent But syntax environment provides better alignment for complex types:

\bigskip

\begin{syntax}
Color2 & ::= & red2 | green2 | blue2 \\
Shape & ::= & circle \ldata \nat \rdata | rectangle \ldata \nat \cross \nat \rdata | polygon \ldata \seq~\nat \rdata
\end{syntax}

\noindent The column alignment makes the structure more readable.

\bigskip

\end{document}